\section{Introduction}
\label{sec:intro}

Randomized experiments are often run in sites to establish evidence of an intervention applied in another site. When site-level context changes the treatment response, a
pool of internally valid experiments identifies the effect at a new site
only in part. \citet{boyarsky2026} formalize this problem. Site context
enters each potential-outcome surface through linear fixed effects, the
observed sites determine those fixed effects only inside the span of
their centered context profiles, and the target effect is reported as an
interval: the smallest and largest values consistent with the source
experiments and with restrictions the analyst is willing to state. A wide interval means that the data and the stated restrictions do not
settle the target conclusion. An infinite interval means that some
direction of extrapolation is not controlled at all.

An analyst facing a wide interval has two options. She can impose stronger restrictions, which tightens the set at the cost
of assumptions that the data cannot check. She can also collect more
evidence by running an experiment in a new site. The companion paper studies the first option.
This paper studies the second. The question is concrete because site
context is observable before any experiment is run. A funder choosing
among candidate countries already knows each country's interest rates,
loan sizes, and program rules. What she does not know is what an
experiment there would find. We ask which candidate site, if its
experiment were run, would most strengthen the external validity of the
pooled evidence, and how much of that answer is available before the
experiment.

We use the microcredit experiments harmonized by
\citet{meager2019understanding} as a running example. Seven countries ran nominally similar
expansions of credit access under very different lending terms. The
companion paper's case study holds out one country at a time and asks
what the other six imply for it. Its leave-one-out exercise shows the reverse of our question. Dropping
any one of Bosnia's six source sites leaves Bosnia's interval unsigned.
That finding also says that the sixth experiment was worth running. We ask whether an analyst who had only five
source sites could have known this in advance, from lending terms alone, and
could have chosen the right country to add.

The paper has three main findings. Each follows from the identification
geometry of the companion paper.

First, site selection is a question about the affine hull of the site
profiles. In the companion framework the observed sites identify the
fixed effects only inside the span of their centered profiles, and we
show that the target effect is point identified exactly when the target
profile lies in the affine hull of the observed profiles
(\cref{sec:geometry}). A candidate site changes what is identified only if its profile lies
outside that hull. Whether it does is checkable at design time, before
any data are collected. How far outside it lies does not matter for
identification, because any experiment along the new direction pins down
the same coordinate of the fixed effects. Distance affects precision but
not the identified set. It also affects robustness, since a farther site
is less exposed to a site-level conflict with the model.

Second, the value of a candidate experiment is computable in advance up
to one scalar per treatment arm (\cref{sec:selection}). Under the
maintained model, adding site $k$ replaces the feasible set
$\feasset$ by $\feasset$ intersected with two equality constraints,
$\newdir{k}^\top\beta_a=\reveal{a}{k}$ for $a\in\{0,1\}$, where the
direction $\newdir{k}$ is known from the candidate's context and the
scalars $\reveal{a}{k}$ are what the experiment reveals. The whole map
from those scalars to the post-addition bounds can be solved before the
experiment. Forecasting the scalars from the existing pool gives a predicted
interval, and the dual variable of the added constraint gives the effect
of a forecast error on the bounds. This is the design rule we apply in
the simulations and the case study.

Third, under the maintained model a new site can never hurt. Every width
weakly shrinks, at every target (\cref{sec:monotonicity}). This may seem
obvious, but it fails once the model is wrong. A site whose outcomes
conflict with the shared response model moves the pooled moment
conditions to a compromise between the sites. The identified slice then
translates, and the width can grow or shrink, or the program can become
infeasible. We
give an example with two source sites where one conflicting addition
doubles the width. The effect of the conflict on correctness is more important than its
effect on the width. The translated set stops covering the true effect
exactly at the targets the pool point identifies. The screens we
propose, a rank test on profiles and a compatibility check on outcomes,
separate the interior additions that refine the set from those that
corrupt it. For a site outside the hull no such check exists. We show
that a constant conflict at such a site is exactly what a different
coefficient vector would produce under the model, so the source data
cannot reveal it.

We verify each claim with the estimator the companion paper recommends,
not only with population calculations. The simulations in
\cref{sec:simulation} generate multi-site data from the fixed-effects
model, estimate the bounds with the cross-fitted rank consistent
estimator, and show that the estimated selection recovers the
population-optimal candidate in every replicate, that the design-time
predictions match the realized widths, and that the conflict mechanism
behaves as the theory says. The case study in \cref{sec:case-study}
runs the same exercise on the microcredit data. For each five-country
source pool, the rule scores the missing country from its lending terms,
and the realized experiment confirms the prediction.

\paragraph{Related work.}
The site-selection question is the forward version of the leave-one-out
exercise in the case study of \citet{boyarsky2026}, which asks which observed
sites a reported bound depends on. We ask which unobserved site a future
bound would benefit from, and we use the same three categories of what a site can do: expand
identification, add precision, or break feasibility. The problem is also related to classical optimal experimental design \citep{atkinson2007optimum,pukelsheim2006optimal}. There the
design variable is an allocation of observations inside one population
and the loss is a variance. Here the design variable is a site-level
context profile and the loss is the width of a partial-identification
interval, which variance reduction alone cannot shrink. Submodular
coverage arguments for sensor placement and subset selection
\citep{nemhauser1978analysis} appear when several sites are added at
once (\cref{sec:discussion}). Distributionally robust formulations of
worst-case effects \citep{jeong2020assessing} and structural
transportability conditions \citep{bareinboim2016causal} address what
can be learned from fixed data. Our contribution is to treat the data
that do not yet exist as the design object.
